\documentclass[10pt,a4paper]{article}
\usepackage[margin=1.35cm]{geometry}
\usepackage{amsmath,amssymb}
\usepackage{booktabs}
\usepackage{array}
\usepackage{xcolor}
\usepackage{enumitem}
\usepackage[hidelinks]{hyperref}
\usepackage{caption}
\usepackage{multicol}
\usepackage{titlesec}

\titlespacing*{\section}{0pt}{6pt}{3pt}
\titlespacing*{\subsection}{0pt}{4pt}{2pt}
\titleformat{\section}{\normalfont\large\bfseries}{\thesection}{0.6em}{}
\titleformat{\subsection}{\normalfont\normalsize\bfseries}{\thesubsection}{0.6em}{}
\setlist{nosep,leftmargin=1.2em}
\captionsetup{font=small,skip=2pt}
\pagestyle{empty}
\setlength{\parskip}{2pt}
\setlength{\tabcolsep}{4pt}
\renewcommand{\arraystretch}{0.95}

\newcommand{\F}[1]{\textcolor{blue!55!black}{\footnotesize[#1]}}

\begin{document}

\begin{center}
{\large\textbf{CS4.406 Assignment 1 --- Lexical \& Semantic Retrieval: Design Note}}\\[2pt]
{\small Noel Alex Jacob (2025201085) \quad·\quad
\href{https://github.com/N03L-22/IRE_A1_C1}{\texttt{github.com/N03L-22/IRE\_A1\_C1}} \quad·\quad
95 tests \quad·\quad \texttt{make data \&\& make store \&\& make eval}}
\end{center}
\vspace{-4pt}

\section{What I built}

Raw archives $\rightarrow$ per-dataset readers $\rightarrow$ \textbf{one unified schema} $\rightarrow$ temporal split
$\rightarrow$ parquet feature store (42\,s rebuild, both datasets). Two retrievers behind one
\texttt{Retriever} protocol --- BM25 over title+abstract, and MiniLM embeddings with a brute-force
exact index --- plus recency, popularity, random, RRF fusion and a popularity-prior blend. One
\texttt{evaluate()} scores all of them, so Q4.5 holds by construction: there is nowhere else a metric
can be computed.

\textbf{Two evaluation regimes are kept strictly apart}, because conflating them produces numbers
that look fine and mean nothing. \emph{Corpus regime}: recall@K over all 21K--65K articles ---
\emph{did the clicked article survive the cut?} \emph{Slate regime}: AUC/MRR/nDCG over the
$\sim$11--40 candidates actually shown --- \emph{did you rank it above the others?} Every result row
carries a \texttt{regime} column.

\section{The constraint that drove every design choice}

Every richer feature EB-NeRD offers is one MIND lacks: body text, click timestamps, publish time,
session ids (0\% on MIND), provided embeddings. Using any of them for a headline number silently
converts a \emph{dataset} comparison into a \emph{feature-availability} comparison.

\textbf{Rule adopted: the shared path uses only the intersection} --- title, abstract, category,
click history. Everything one-sided becomes a labelled single-dataset ablation. This one rule
resolves the encoder choice, recency weighting, body text and session context together.

\section{Key decisions and rejected alternatives}

\begin{center}\footnotesize
\begin{tabular}{@{}p{2.4cm}p{5.4cm}p{9.3cm}@{}}
\toprule
\textbf{Decision} & \textbf{Chosen} & \textbf{Rejected, and why} \\
\midrule
Corpus & Union of MIND's two \texttt{news.tsv} (65,238) & Per-split: 23.3\% of dev clicks are on train-absent articles, an \textbf{invisible 0.767 recall ceiling} \F{F2,F20} \\
Split & Hold out official test, carve val from train tail & Literal ``7 days test'': MIND's \emph{entire} labelled span is 6.0 days \F{F27} \\
Lexical fields & Title + abstract (Q2.1 mandates) & Title+body --- MIND has no body, breaking comparability \F{F10} \\
Encoder & \textbf{MiniLM 384-d}, mean-pooled, L2-normalised & \textbf{XLM-R fails the Danish probe} (below); \texttt{[CLS]} pooling; per-language stemming \\
Index & Brute force (exact; Q3.2 permits it) & HNSW --- both corpora are $<$200K, where exact costs 190--890\,ms and is the recall \emph{ceiling} ANN is measured against \\
User vector & Log positional decay, $n{=}20$, conditional pooling & Exponential decay collapses the query to 2--3 clicks; plain mean blurs multi-interest users \\
Session context & \textbf{Dropped}, deferred to C-2 & 0\% coverage on MIND; EB-NeRD sessions average \textbf{1.9 impressions} \F{F29} \\
Coverage CI & Point estimate, marked \texttt{(no CI)} & Every bootstrap scheme is biased low --- see \S\ref{sec:ci} \\
\bottomrule
\end{tabular}
\end{center}

\section{Observations}

\subsection{Recency dominates; text matching is secondary \F{F16,F21}}

On EB-NeRD, \textbf{94.3\% of clicks are on articles under 24\,h old} while only $\sim$1.1\% of the
corpus is that fresh. A retriever that \emph{ignores the user entirely} and returns the newest
articles scores \textbf{recall@50 = 0.8986}, against BM25's 0.0043 over the full corpus. Restricting
either retriever to a 24\,h window moves it from $\sim$0.004 to $\sim$0.24 --- a \textbf{33$\times$
swing from the candidate pool, not the scoring function}.

Candidate generation on news is primarily a \emph{freshness} problem. This is legitimate, not a leak:
publish time is known at serving time; ranking by \emph{future} engagement would not be.

\subsection{Lexical vs. semantic: the answer flips between datasets \F{Q3.5}}

\begin{center}\footnotesize
\begin{tabular}{@{}lll@{}}
\toprule
\textbf{recall@100}, $n{=}4{,}000$ & \textbf{EB-NeRD (Danish)} & \textbf{MIND (English)} \\
\midrule
Best overall & recency \textbf{0.9434} [0.9361, 0.9505] & popularity \textbf{0.0688} [0.0615, 0.0763] \\
Lexical & bm25+24h \textbf{0.2375} [0.2244, 0.2490] & bm25 0.0142 [0.0108, 0.0172] \\
Semantic & sem+24h 0.2307 [0.2175, 0.2432] & sem \textbf{0.0163} [0.0130, 0.0199] \\
RRF fusion & 0.0070 --- \emph{no gain} & \textbf{0.0181} [0.0146, 0.0218] --- \emph{best retriever} \\
\bottomrule
\end{tabular}
\end{center}

\textbf{The point estimates favour semantic on MIND and lexical on EB-NeRD --- but a paired
bootstrap on the difference finds \emph{neither gap significant}}: MIND fusion vs BM25 nDCG@10
$+0.0033$ [$-0.0029$, $+0.0105$] (they differ on 55 of 800 impressions); EB-NeRD semantic vs BM25
$-0.0055$ [$-0.0277$, $+0.0155$] (differing on 470 of 800 --- ample signal, no detectable gap). At
$n=4{,}000$ the two families are indistinguishable and the apparent flip is not established.

\textbf{The leaderboard disagrees, and that is the finding.} Submitting fusion for MIND scored
\textbf{AUC 0.5934 against BM25's 0.5568} --- a real $+0.0366$ on 2.37M impressions, an effect the
offline harness had no power to see. An offline proxy this size can rule out large differences; it
cannot rank retrievers.

\textbf{The slice tells a sharper story.} EB-NeRD recall@100, cold defined as history $\le$ 97
(q0.25 of observed; median 235):

\begin{center}\footnotesize
\begin{tabular}{@{}lll@{}}
\toprule
 & \textbf{cold} ($n{=}1{,}004$) & \textbf{warm} ($n{=}2{,}996$) \\
\midrule
semantic + 24h & \textbf{0.2468} [0.2198, 0.2742] & 0.2253 [0.2104, 0.2407] \\
bm25 + 24h     & 0.2216 [0.1982, 0.2490] & \textbf{0.2428} [0.2276, 0.2588] \\
\bottomrule
\end{tabular}
\end{center}

\textbf{The ordering crosses over}: semantic wins cold users, lexical wins warm ones. With few
clicks, word overlap has too few terms to work with while embeddings still place a sparse history
somewhere meaningful. CIs overlap, so this is suggestive, not established --- but the crossover is
the shape, and it is what Q3.5 asks for.

\subsection{XLM-RoBERTa cannot separate related from unrelated Danish \F{F37}}

The brief names BERT/XLM-RoBERTa. Before trusting either, a probe embedded 5 obviously-related
Danish headline pairs and 5 obviously-unrelated ones:

\begin{center}\footnotesize
\begin{tabular}{@{}lrrrl@{}}
\toprule
Encoder & related & unrelated & margin & verdict \\
\midrule
\texttt{xlm-roberta-base} (768-d) & 0.9972 & 0.9954 & \textbf{+0.0018} & \textbf{OVERLAPS} \\
MiniLM (384-d) & 0.6523 & 0.0253 & \textbf{+0.6271} & SEPARATES \\
\bottomrule
\end{tabular}
\end{center}

XLM-R rates \emph{``Brøndby beat FCK''} and \emph{``a new apple cake recipe''} as 0.995 similar. It
is trained for masked-language modelling, not to make cosine similarity meaningful, and its vectors
collapse into a narrow cone. \textbf{A retriever built on it returns arbitrary articles while
producing perfectly plausible metrics} --- the failure mode this probe exists to catch. MiniLM
achieves a 348$\times$ larger margin from \emph{half} the dimensions. Reported as a measured failure
rather than a silent substitution.

\subsection{Slate metrics barely discriminate --- always show the random row \F{F25}}

EB-NeRD slates average 11--12 items with exactly one click, so \textbf{random ranking scores
nDCG@10 = 0.4367 and AUC = 0.4933}. The best retriever reaches 0.4750 / 0.5407. Every AUC sits
within noise of 0.5. Any slate-regime table without a random baseline row is misleading, and
recall@K in the corpus regime is the only metric that clearly separates retrievers here.

\subsection{Conditional pooling earns its complexity on one dataset only}

Pooling strategy counts over 4,000 impressions: EB-NeRD \texttt{mean} 3,991 / \texttt{recent\_half}
9; MIND \texttt{mean} 2,117 / \texttt{recent\_half} 1,462 / \texttt{max\_pool} 270. \textbf{46\% of
MIND users have interests too scattered for a centroid, against 0.2\% on EB-NeRD} --- consistent
with MIND's shorter, more varied histories (median 19 vs 92 clicks).

\subsection{One metric legitimately has no confidence interval}\label{sec:ci}

Q4.4 demands a bootstrap CI on every metric. \textbf{Coverage cannot honestly have one.} It counts
\emph{distinct} articles, so it grows with sample size; resampling $n$ units with replacement makes
$\sim$37\% duplicates, so every resample sees fewer unique articles. The first attempt produced
$0.9783\ [0.9035, 0.9235]$ --- \textbf{the point estimate outside its own interval}. Subsampling
without replacement fails at every ratio from 0.5 to 0.99. Reported as a point estimate marked
\texttt{(no CI)}: saying so beats shipping a plausible-looking number that is biased by construction.

\subsection{Embedding width: 256-d beats 384-d, and 128-d costs nothing \F{F47}}

Dimension is the one semantic knob with a direct resource cost --- memory, index size and search
time all scale linearly in it. Truncating the cached 384-d MiniLM vectors and re-normalising
isolates it exactly: same model, same training, only the width changes.

\begin{center}\footnotesize
\begin{tabular}{@{}lrrl@{}}
\toprule
Dim & Vectors & recall@50 & Paired difference vs 384-d \\
\midrule
384 & 31.9 MB & 0.2325 & --- \\
\textbf{256} & \textbf{21.2 MB} & \textbf{0.2500} & \textbf{$+0.0175$ [$+0.0025$, $+0.0338$] SIGNIFICANT} \\
\textbf{128} & \textbf{10.6 MB} & 0.2437 & $+0.0112$ [$-0.0112$, $+0.0338$] --- no worse \\
64 & 5.3 MB & 0.2175 & $-0.0150$ --- degrades \\
\bottomrule
\end{tabular}
\end{center}

\textbf{256-d is significantly better than full width at 34\% less memory}, and 128-d is
indistinguishable from it at a third. The tail dimensions of a non-Matryoshka embedding carry mostly
noise for this task; dropping them removes variance without removing signal --- the mechanism that
makes PCA often improve retrieval rather than merely compress it. Combined with \F{F37}, where
XLM-R's 768-d lost to MiniLM's 384-d by 348$\times$ on the Danish probe: \textbf{the training
objective and the noise floor decide quality, the width does not}.

\subsection{Query length: an inverted U, and still not significant \F{F51}}

Swept \texttt{last\_n} \emph{inside} the 24\,h window, the regime we ship, rather than over the full
corpus where everything scores under 0.008. recall@50: 0.2288 (5), 0.2313 (10), 0.2475 (15),
\textbf{0.2531 (20)}, 0.2456 (30), 0.2338 (50) --- a clear peak at 20. Paired against the shipped 15
it is $+0.0056$ [$-0.0156$, $+0.0250$], \textbf{not significant}, and neither is any other value.
The shape matches the full-corpus sweep, so the window does not change this conclusion even though
window and $K$ do interact elsewhere. We ship 15 and report that 20 is a candidate, not an
improvement.

\subsection{Multi-query retrieval: built, measured, rejected \F{F48}}

The stated weakness of a single user vector is that someone who reads football \emph{and} recipes
gets a centroid between the two, matching neither. Clustering the history and retrieving per centroid
should fix that --- and \F{F40} found \textbf{46\% of MIND users} falling back from the mean, so the
population exists.

Measured: \textbf{not significant on either dataset, at 19--36$\times$ the query cost} (MIND 57.5\,s
vs 3.0\,s; paired $+0.0009$ [$-0.0055$, $+0.0077$]). The clustering fired properly --- 1,518 of 1,564
MIND users received three centroids --- yet the two variants differ on only \textbf{16 of 800}
impressions.

\textbf{So the blended centroid is not as useless as the theory assumes.} It still ranks each
interest's articles highly enough to reach a top-50, because \emph{recall@K at K{=}50 is a forgiving
target}. The blur would matter for precision at rank 1 --- a re-ranker's problem, not candidate
generation's. Keeping the single vector was the right call, and is now evidence rather than a guess.

\subsection{Overlapping intervals are the wrong test \F{F46}}

``Do the CIs overlap?'' is a conservative approximation. Both retrievers are scored on the
\emph{same} impressions, so their shared per-impression noise cancels when you subtract --- and a
paired bootstrap on the \emph{difference} is both the correct test and a much tighter one.

It changed two conclusions. The 384-vs-256 intervals above overlap heavily ([0.2025, 0.2613] against
[0.2206, 0.2806]) and the overlap heuristic would have called it undecided; the paired test finds it
significant, so we would have shipped the larger, worse vector. Conversely an earlier draft of this
note claimed semantic ``edges ahead'' of lexical on EB-NeRD from the point estimates; paired, the
difference is $-0.0055$ [$-0.0277$, $+0.0155$] --- the sign reverses and nothing is established.

\textbf{Every null result now reports how many impressions the two systems actually differ on},
because that separates \emph{no effect} from \emph{no power}: the dedup ablation differed on
\textbf{4 of 800} (unmeasurable by construction), while EB-NeRD semantic-vs-lexical differed on
\textbf{470 of 800} --- abundant signal, genuinely equal.

\section{Anti-gaming (Q9)}

\textbf{Leakage.} History is truncated to clicks strictly before the impression time.
\texttt{test\_no\_leakage.py} asserts this on the built store \emph{and injects a deliberate
post-boundary click to confirm the checker fails} --- a test that passes both before and after the
boundary breaks proves nothing. \emph{Honest limit:} verifiable only where per-click timestamps
exist, i.e.\ EB-NeRD. MIND's history is untimestamped, so there the boundary rests on the authors'
construction; the store records \texttt{history\_boundary\_verifiable: false} rather than implying
otherwise.

\textbf{Serving-unavailable features.} EB-NeRD's training rows carry \texttt{read\_time},
\texttt{scroll\_percentage} and \texttt{next\_read\_time}/\texttt{next\_scroll\_percentage} --- all
recorded \emph{after} the click. Predicting a click from its own consequences is circular. The
organisers settled the boundary: those four columns are exactly what is deleted from the test set.
They are excluded in \texttt{clean.py} with a comment naming why. Our retrievers never touch them,
so the honest with/without pair is the cold-start fallback flag, reported as a \texttt{features}
column on every result row.

\section{Where it breaks at 10$\times$}

\begin{itemize}
\item \textbf{The join breaks before the model does \F{F15}.} EB-NeRD large is 52$\times$ small's
impressions and its history parquet exceeds its behaviours file. The data layer gives first.
\item \textbf{The prediction path needed a different algorithm, not a bigger machine \F{F32}.}
Full-corpus retrieval per impression ran at 162 impressions/s ($\sim$4\,h/file) because it discarded
99\% of each result. Scoring the slate directly: \textbf{2,570/s, 16$\times$}.
\item \textbf{Parallelism was memory-bound, and the fix was representation \F{F36}.} A worker peaked
at 15.1\,GB, of which the index was only 1.8\,GB --- the rest was 116.8M click ids as Python
\texttt{str}. Storing them as \texttt{array('i')} indices: \textbf{9.25\,GB, and 45\% faster}.
\item \textbf{Swap does not slow the merge, it stops it \F{F38}.} All 51 shards wrote at 6,426
lines/s, then the merge produced \emph{zero} output for 20+ minutes with 24\,GB swapped. The same
merge with RAM free: \textbf{7 seconds}. Scoring streams and degrades gracefully; the merge hits a
13.3M-element set at random, where swap does not degrade --- it stalls.
\item \textbf{ANN: the trade-off only becomes real at 10$\times$ \F{F49/F50}.} On \emph{clustered}
vectors HNSW is 62--69$\times$ faster than exact at \textbf{0.99+ recall} --- an earlier benchmark
using uniform-random vectors measured only 0.45--0.77 and was a synthetic-data artefact, since
random vectors in 384 dimensions are near-orthogonal and give a proximity graph nothing to exploit.
At 1M vectors the speed-up grows to 160$\times$ but \textbf{recall falls to 0.8332}: exact search
becomes unusable (4.6\,s/query) exactly as the approximation starts costing 17\% of the answer.
\textbf{That is where the pipeline breaks} --- not on memory or the model, but on the ANN index's
recall/latency trade-off ceasing to be free. Build time also grows super-linearly, 4\,s
$\rightarrow$ 156\,s.
\end{itemize}

\section{Leaderboards (Q5) and limitations}

MIND submitted: \textbf{AUC 0.5568, MRR 0.2646, nDCG@5 0.2778, nDCG@10 0.3331}. EB-NeRD file
generated and validated (13,336,711 lines, 100\% valid permutations). \emph{Screenshots overleaf.}

\textbf{Calibration \F{F44}.} The challenge paper (arXiv 2409.20483) publishes reference points for
EB-NeRD: winners at 89.24 / 88.15 / 87.07 AUC, a \emph{most-clicks} editorial baseline at
\textbf{59.70}, read-time 59.49, in-view-rate 54.50, random 49.98. Our MIND fusion at 0.5934 sits
essentially at the most-clicks baseline --- an analogy across datasets, not a like-for-like, but a
useful order of magnitude. All three winners used GBDT ensembles or transformers in multi-stage
pipelines with engineered temporal features, i.e.\ a \textbf{re-ranker}: the 0.59$\rightarrow$0.87
gap is a different class of system, not a better-tuned candidate generator. Notably their ablations
independently identified \emph{popularity} and \emph{article timeliness} as the load-bearing
signals --- the same two levers this component found in \F{F16/F41} and \F{F25/F39}.

\textbf{Limitations, stated plainly.} (i)~The offline harness was under-powered: at $n{=}800$ it
reported MIND AUC 0.4981 [0.4776, 0.5190] while the leaderboard scored \textbf{0.5568} ---
\emph{outside} the interval \F{F34}. Now $n{=}4{,}000$+, but several reported differences still have
overlapping CIs and are labelled as such. (ii)~Abstracts are 75\% of the lexical index and buy
\textbf{+0.011 recall against a $\pm$0.030 CI} --- kept because Q2.1 mandates it, reported as the
null result it is \F{F28}. (iii)~HNSW figures come from \emph{random} vectors (near-orthogonal,
worst case) and are a pessimistic bound. (iv)~MIND's test window is a single day, so one day's
topical anomaly moves every MIND number.

\end{document}
